\documentclass[10pt,a4paper]{amsart}
\usepackage[margin=19mm]{geometry}
\usepackage{amsmath,amssymb,mathtools}
\usepackage[scr=rsfs,cal=euler]{mathalfa}
\usepackage{xfrac}
\usepackage[unicode,hidelinks,bookmarks=false]{hyperref}
\newcommand{\ie}{i.e.\ }
\newcommand{\cf}{cf.\ }
\newcommand{\resp}{respectively\ }
\newcommand{\Subsection}{\subsection}
\providecommand{\nobreakdash}{\nobreak\hbox{-}}
\setlength{\emergencystretch}{2em}
\multlinegap0pt
\usepackage{bm}
\usepackage{bbm}
\usepackage{dsfont}
\usepackage{xspace}
\usepackage{cancel}
\usepackage{mathtools}
\usepackage{cleveref}
\usepackage{amsthm}
\makeatletter
\@ifundefined{thm}{\newtheorem{thm}{Theorem}}{}
\@ifundefined{cor}{\newtheorem{cor}[thm]{Corollary}}{}
\@ifundefined{lem}{\newtheorem{lem}[thm]{Lemma}}{}
\makeatother
\makeatletter
\newcommand{\cC}{\mathcal{C}}
\newcommand{\cN}{\mathcal{N}}
\@ifundefined{claimproof}{\newenvironment{claimproof}[1][\claimproofname]{\begin{proof}[#1]\renewcommand*{\qedsymbol}{\(\bullet\)}}{\end{proof}}}{}
\@ifundefined{symenum}{\newenvironment{symenum}
 {\enumerate[label=\normalfont{(\noexpand\thisenumsymbol)},align=parleft,labelindent=0pt,itemindent=0pt,labelsep=35pt,leftmargin=*]}}{}
\makeatother
\title[A hierarchy of edge-weight symmetries in perfect matchings]{A hierarchy of edge-weight symmetries in perfect matchings}
\author{Krist\'of B\'erczi, Viktor Csapl\'ar, Yutaro Yamaguchi}
\date{Source: arXiv:2604.20009v1 (CC BY 4.0)}
\begin{document}
\maketitle

\noindent\emph{Source and licence.} A hierarchy of edge-weight symmetries in perfect matchings. Version arXiv:2604.20009v1 (CC BY 4.0), distributed under the Creative Commons Attribution 4.0 International licence, as recorded in the arXiv licence metadata for this exact version and shown on the abstract page https://arxiv.org/abs/2604.20009v1. Full rights notice: licenses/arxiv-2604.20009-CC-BY-4.0.txt.\par\smallskip

\noindent\textit{Editorial scope.} Counted unit: the Thm whose source label is \texttt{thm:ear\_decomposition\_main} in the article (source0.tex lines 648--650, source label \texttt{thm:ear\_decomposition\_main}), delivered together with the article's own complete original proof of it (source0.tex lines 651--689, one \texttt{proof} environment of 39 lines). No mathematics is written, repaired or re-derived: statement and proof are the article's own text line for line. Required context retained uncounted: thm:even\_odd (lines 407--409); thm:whitney (lines 429--431); lem:space (lines 590--592); cor:jgnofg (lines 616--618); lem:C\_subset\_of\_N\_perp (lines 639--641). Every definition, display and label of those lines is retained unchanged. Omitted from this package and not redistributed: 1--406, 410--428, 432--589, 593--615, 619--638, 642--647, 690--896. Nothing inside a retained excerpt is omitted or altered: every retained body is delivered exactly as the article's lines are written, so no omission receipt of any kind accompanies this package. Citations: the retained excerpts carry no citation command at all, so no bibliography entry of the article is transcribed and no reference is redistributed. Reference adaptation: none was needed and none was made; every reference inside the delivered unit resolves inside it, so no unresolved reference or citation is produced. Printed slips kept literal: no printed word, symbol, hypothesis or number of the counted statement or of its proof was altered. Layout and class: the article's own class is \texttt{article} with packages including fullpage, inputenc, amsmath, amsthm, amssymb, nicefrac, amsfonts, thmtools, mathtools, leftindex, leftidx, caption, subcaption, thm-restate; the wrapper is the lane's pinned \texttt{amsart} editorial wrapper at 10pt on A4, which restates the theorem-like environments the retained text needs and adds the article's own macro definitions that the retained text needs (\texttt{\textbackslash cC}, \texttt{\textbackslash cN}), with extra wrapper packages (bm, bbm, dsfont, xspace, cancel, mathtools, cleveref). The delivered numbering is the unit's own. Assets: no retained span contains a figure environment, a graphics-inclusion command, a PDF-page inclusion, a table or a TikZ picture; no image file is copied into this package, the package declares no asset dependency and no image file is redistributed with it. Licence: version arXiv:2604.20009v1 of this article is distributed under the Creative Commons Attribution 4.0 International licence, as recorded in the arXiv licence metadata for this exact version and shown on the abstract page https://arxiv.org/abs/2604.20009v1. A full rights notice is delivered at licenses/arxiv-2604.20009-CC-BY-4.0.txt. Characters: every retained excerpt is pure ASCII, so the delivered engine, XeLaTeX with its default Latin Modern text font, reports no missing character.
\begin{thm}[Little, Vince]\label{thm:even_odd}
    Let $C$ be an odd cycle in a $2$-connected graph $G$, and let $e \in E(G) \setminus E(C)$. Then $G$ contains both an even and an odd cycle through $e$ that also contains an edge of $C$.
\end{thm}

\begin{thm}[Robbins, Whitney]\label{thm:whitney}
   A graph admits an ear decomposition if and only if it is 2-edge-connected, and an open ear decomposition if and only if it is 2-connected. Furthermore, for any 2-edge-connected graph $G$, there exists an ear decomposition starting from an arbitrary block $B_0$ and any cycle $C\subseteq B_0$ as the first ear, such that a closed ear is introduced if and only if the decomposition enters a new block.
\end{thm}

\begin{lem}\label{lem:space}
    Let $G = (V, E)$ be a connected graph. Then $\dim(\mathcal{N}_G) = |V| - 1$ if $G$ is bipartite and $\dim(\mathcal{N}_G) = |V|$ otherwise.
\end{lem}

\begin{cor}\label{cor:jgnofg}
Let $G$ be a loopless graph. Then $\dim(\cN_G) = |V| - p(G)$, where $p(G)$ denotes the number of bipartite components of $G$.
\end{cor}

\begin{lem}\label{lem:C_subset_of_N_perp}
    For any graph $G$, we have $\cC_G \subseteq \cN_G^{\perp}$.
\end{lem}

\begin{thm}\label{thm:ear_decomposition_main}
    Let $G$ be a graph whose block decomposition contains at most one non-bipartite block. Then $\cC_G = \cN_G^{\perp}$.
\end{thm}

\begin{proof}
    By \cref{lem:C_subset_of_N_perp}, it is enough to show that $\dim(\mathcal{C}_G) + \dim(\mathcal{N}_G) = |E(G)|$. We may assume that the graph is 2-edge-connected, since if the equality holds for every 2-edge-connected component, then adding each cut edge increases the number of edges by one and decreases the number of bipartite components by one, while the dimension of $\mathcal{C}_G$ remains unchanged. Hence, the required equality follows by Corollary~\ref{cor:jgnofg}.
    
    Therefore, assume that $G$ is 2-edge-connected. We prove this by induction on the number of ears in a carefully chosen ear decomposition given by Theorem~\ref{thm:whitney}: the decomposition is selected so that, if $G$ contains a non-bipartite block, the construction starts within such a block and begins with an odd cycle in that block.

    The base case is $G$ being a cycle. By \cref{lem:space}, $\dim(\cN_G) = |V(G)|-1$ or $|V(G)|$, depending on whether this cycle has even or odd length, respectively. If $G$ is an even cycle, then $\cC_G$ is spanned by its alternating characteristic vector, so $\dim(\cC_G)=1$. Otherwise, there are no even cycles, hence $\dim(\cC_G)=0$. In both cases, the equality $\dim(\cC_G) + \dim(\cN_G) = |E(G)|$ holds, since $|E(G)|=|V(G)|$.
    
    Now assume that $G_{i+1}$ is obtained from $G_i$ by adding one ear, and that the statement holds for $G_i$, i.e. $\cC_{G_i}^{\perp} = \cN_{G_i}$. We first observe that any even cycle of $G_i$ is also an even cycle of $G_{i+1}$, hence $\cC_{G_i} \subseteq \cC_{G_{i+1}}$. Let the added ear introduce $\ell$ new nodes, so that
    $|V(G_{i+1})| = |V(G_i)| + \ell$ and $|E(G_{i+1})| = |E(G_i)| + \ell + 1$. We distinguish two cases depending on bipartiteness.
    \medskip
    
    \noindent\textbf{Case 1:} $G_i$ is bipartite.
    
    If $G_i$ is bipartite, then $G$ is bipartite, since otherwise the decomposition would start with an odd cycle in a non-bipartite block. Hence $G_{i+1}$ is bipartite. By Lemma~\ref{lem:space}, we have $\dim(\cN_{G_{i+1}}) = |V(G_{i+1})| - 1 = |V(G_i)| + \ell - 1$. Since $G_i$ is also bipartite, we similarly have $\dim(\cN_{G_i}) = |V(G_i)| - 1$, which by the induction hypothesis means $\dim(\cC_{G_i}) = |E(G_i)| - |V(G_i)| + 1$. Since $\cC_{G_i} \subseteq \cC_{G_{i+1}} \subseteq \cN_{G_{i+1}}^{\perp}$, we can bound the dimension:
    \begin{align*}
    \dim(\cC_{G_{i+1}}) 
    &\le |E(G_{i+1})| - \dim(\cN_{G_{i+1}})\\
    &= (|E(G_i)| + \ell + 1) - (|V(G_i)| + \ell - 1)\\
    &= |E(G_i)| - |V(G_i)| + 2\\
    &= \dim(\cC_{G_i}) + 1.
    \end{align*}
    We show that $\cC_{G_{i+1}} \neq \cC_{G_i}$. For this, it is enough to find an even cycle in $G_{i+1}$ that uses the new ear. If the ear is closed, then it is itself an even cycle. Otherwise, let its endpoints be $x,y$. Since $G_{i+1}$ remains bipartite, every $xy$-path in $G_i$ has the same parity, hence combining the ear with any $xy$-path yields an even cycle containing the ear.
    \medskip

    \noindent\textbf{Case 2:} $G_i$ is non-bipartite.

    By Lemma~\ref{lem:space}, we have $\dim(\cN_{G_{i+1}}) = |V(G_{i+1})| = |V(G_i)| + \ell$. Since $G_i$ is also non-bipartite, we similarly have $\dim(\cN_{G_i}) = |V(G_i)|$, which by the induction hypothesis means $\dim(\cC_{G_i}) = |E(G_i)| - |V(G_i)|$. Since $\cC_{G_i} \subseteq \cC_{G_{i+1}} \subseteq \cN_{G_{i+1}}^{\perp}$, we can bound the dimension:
    \begin{align*}
    \dim(\cC_{G_{i+1}}) 
    &\le |E(G_{i+1})| - \dim(\cN_{G_{i+1}})\\
    &= (|E(G_i)| + \ell + 1) - (|V(G_i)| + \ell)\\
    &= |E(G_i)| - |V(G_i)| + 1\\
    &= \dim(\cC_{G_i}) + 1.
    \end{align*}
    We show that $\cC_{G_{i+1}} \neq \cC_{G_i}$. For this, it is enough to find an even cycle in $G_{i+1}$ that uses the new ear. If the ear is closed, then it is itself an even cycle as, by the choice of the decomposition, it is part of a bipartite block of $G$. Otherwise, let its endpoints be $x,y$. By Theorem~\ref{thm:even_odd}, there exists an even cycle in $G_{i+1}$ that uses the new ear. Thus $\dim(\cC_{G_{i+1}}) = \dim(\cC_{G_i}) + 1$. Therefore, $\dim(\cC_{G_{i+1}}) + \dim(\cN_{G_{i+1}}) = (|E(G_i)| - |V(G_i)| + 2) + (|V(G_i)| + \ell - 1) = |E(G_{i+1})|$.
    \medskip
    
    In all cases, the inductive step preserves $\dim(\cC_{G_{i+1}}) + \dim(\cN_{G_{i+1}}) = |E(G_{i+1})|$, which completes the proof.
\end{proof}

\end{document}
